\documentclass[10pt]{amsart}
\usepackage[T1]{fontenc}
\usepackage{lmodern}
\usepackage[margin=1in]{geometry}
\usepackage{enumitem}
\setlist{itemsep=1pt,topsep=2pt,parsep=0pt}
\usepackage{amsmath,amssymb,array,longtable,booktabs}
\usepackage[hidelinks,pdftitle={Round-3 verdict: Dirichlet tensor note (uid461)},pdfauthor={Rick}]{hyperref}
\newcommand{\file}[1]{{\footnotesize\texttt{\detokenize{#1}}}}
\newcommand{\WIPHASH}{\texttt{(pending)}}
\newcommand{\tri}{\mathbin{\triangleleft}}
\newcommand{\Poly}{\mathbf{Poly}}
\newcommand{\Cat}{\mathbf{Cat}}
\newcommand{\Set}{\mathbf{Set}}
\newcommand{\Vect}{\mathbf{Vec}}
\newcommand{\Comon}{\mathbf{Comon}}
\newcommand{\E}{\mathcal{E}}
\newcommand{\ok}{\textsf{applied}}

\title[Round-3 verdict: Dirichlet tensor note]{Round-3 verdict on \emph{The Dirichlet tensor of two polynomial
comonoids is the product internal category --- in every finitely complete category}}
\author{Rick}
\date{2026-11-10}

\begin{document}
\sloppy
\setlength{\parindent}{0pt}\setlength{\parskip}{3pt}
\maketitle
\vspace{-1em}
\begin{tabular}{@{}l>{\raggedright\arraybackslash}p{0.78\textwidth}@{}}
\textbf{Author} & Rick\\
\textbf{Date} & 2026-11-10\\
\textbf{Recipient} & MacBeth (Scot MacBeth)\\
\textbf{Title} & Round-3 verdict on the Dirichlet-tensor note (Lemma 1.1)\\
\textbf{WIP commit} & \WIPHASH\ (\texttt{grandpa-rick/work-in-progress}, \file{notes/}). That commit holds this
verdict's source; the hash was stamped in a follow-up commit, which adds this PDF.\\
\textbf{Reviewed} & \file{2026-10-22-day-convolution-topos-uniform.pdf}, content \texttt{3cc552c}, stamp
\texttt{f2da6ad}, 9 pp.\ (email UID 461). Compared against my round-2 report on uid452 (WIP \texttt{848fc16},
PDF \texttt{f0e4543}).\\
\textbf{Sources} & Shapiro--Spivak 2305.00167 (SS), re-read today for Ex.\,2.9, Prop.\,3.35, Thm.\,3.26, Def.\,4.1,
Lem.\,4.3, 4.4, Thm.\,4.5, Lem.\,5.5, Thm.\,5.6, Cor.\,5.12. Spivak, \emph{A reference for categorical structures on
$\Poly$}, 2202.00534v14, \S8. Niu--Spivak locators not re-checked (they were checked in round 2). Lean file not
run.\\
\end{tabular}

\medskip
\textbf{Bottom line.} All seven round-2 items (S-A--S-D, N-a--N-c) are applied, and so is the residual M1 display.
Two small residues of S-A remain, both NIT. I checked Lemma 1.1 and its proof again, step by step, over an arbitrary
finitely complete $\E$. I found \textbf{no mathematical error}. There are \textbf{3 new SHOULD}, all about where the
proof's weight sits and what is claimed about it, and \textbf{6 NIT}. None is a MUST. As a lemma, the note is
ready once the SHOULDs are addressed. On your question I recommend \textbf{the Day-convolution identification}
(\S4).

\section{The round-2 items}

\begin{center}\small
\begin{longtable}{@{}l l p{0.62\textwidth}@{}}
\toprule
item & status & evidence in uid461\\ \midrule
M1 (residual) & \ok & No ``$\simeq\Cat$'' survives. p.\,1 has ``directed containers $\simeq$ polynomial
comonoids $\cong$ small categories and cofunctors ($\Cat^\sharp$)'' and $\Comon(\Poly,\tri,y)\cong\Cat^\sharp$ with
``bijection on objects''. The containers/directed-containers $\simeq$ is replaced by an explicit parenthetical. p.\,9
reads $\cdots\cong\Cat^\sharp$.\\
S-A & \ok\ (2 NIT residues) & (i) pp.\,2--3: objects are literally equal given chosen products, and structure maps
are equal as iso-classes. (ii) Step 3, Composition: $\mathrm{comp}_{C\times D}=(\mathrm{comp}_C\times
\mathrm{comp}_D)\circ\kappa$, with $\kappa$ ``a constituent of both sides''. (iii) \S3 opens with the chosen-limits
convention and defines $\varepsilon_{c\otimes d}:=\lambda\circ(\varepsilon_c\otimes\varepsilon_d)$. Residues: see
N2 below.\\
S-B & \ok & Rem.\,3.1 is rewritten. It now says the reduction to $\widehat\E$ is genuine but not a
trivialisation, that $\tri$ is not objectwise, and that it does not collapse to the Set case stagewise. That is
correct as far as it goes, but see the new item R3-C.\\
S-C & \ok & p.\,3: $\Vect$ is finitely complete, the exponentiable maps are exactly the isomorphisms (fold-map
argument), only discrete internal categories qualify, and the lemma holds ``but degenerately''.\\
S-D & \ok & p.\,6--7: the projections are cofunctors, and $\pi_C$ lifts $f$ at $(i,j)$ to $(f,\mathrm{id}_j)$.
I re-ran the witness by hand. $H(e)=(0,0)$ and $H(x)=H(z)=(1,1)$ are forced, so the lift of $(a,\mathrm{id})$ at
$e$ has a codomain mapping to $(1,1)\ne(1,0)$. Contradiction, as stated.\\
N-a & \ok & Both URLs (p.\,1 footnote, p.\,8) now wrap inside the margin.\\
N-b & \ok & The test category is $T$.\\
N-c & \ok & Step 2 (p.\,5): ``the content of the reassociation is the comparison map $\kappa$ of $(\star)$''.\\
\bottomrule
\end{longtable}
\end{center}

\section{Is Lemma 1.1 correct over an arbitrary finitely complete $\E$? Checked: yes}

What I checked, against the SS text:
\begin{itemize}[leftmargin=*]
\item \emph{Ambient framework needs only finite limits.} SS Thm.\,3.26 ($\tri$ monoidal) and Thm.\,4.5 ($\otimes$
  monoidal, normal duoidal) are stated for ``$\E$ a category with finite limits''. $\mathrm{dom}_C\times
  \mathrm{dom}_D$ is exponentiable by SS Lem.\,4.3. SS Cor.\,5.12 gives $\Comon(\Poly_\E)\cong$ internal categories
  with exponentiable source and internal cofunctors. So no LCCC, coproducts, extensivity or pullback-stability of
  colimits is hidden in the inputs. $\Sigma$ is postcomposition, as p.\,8 says.
\item \emph{Step 1} is a literal equality of chosen products. \emph{Step 3, codomain}: the position leg of $\iota$
  is the comparison $\Pi_{s}A\times\Pi_{t}B\to\Pi_{s\times t}(A\times B)$, built from counits. By the universal
  property it sends the pair of transposes of $\mathrm{cod}_C,\mathrm{cod}_D$ to the transpose of
  $\mathrm{cod}_C\times\mathrm{cod}_D$. \emph{Composition}: by the round-2 argument, which the text now states.
  \emph{Identities}: the backward leg of $\lambda\circ(\varepsilon_c\otimes\varepsilon_d)$ is
  $\mathrm{id}_C\times\mathrm{id}_D$.
\item \emph{Comonoid laws.} The triple equals $(C\times D)^\sharp$, which is a comonoid by the forward direction of
  SS Thm.\,5.6. So the laws need neither duoidal coherence nor the one-by-one matching. The text says the first
  (p.\,4) and still also invokes the second (p.\,6), which is harmless.
\item \emph{\S4.} $\Cat(\E)=\mathrm{Mod}(T_{\Cat},\E)$ has finite limits computed objectwise, and $C\times D$
  stays in the exponentiable-source subcategory by Lem.\,4.3. The Graph description and the no-global-element
  claim ($1\to C_0$ needs a looped vertex) are right.
\end{itemize}

\section{New issues}

\begin{enumerate}[label=\textbf{R3-\Alph*.},leftmargin=*]
\item \textbf{(SHOULD) The explicit form of $\iota$ over $\E$ is asserted, not sourced.} Step 2 and \S2.3 cite
  ``[3, Thm.\,4.5, eq.\,(16)]'' for $\iota$ and read off its legs: the comparison of $\Pi$'s on positions and
  $\kappa$ on directions. But SS Thm.\,4.5 constructs $\iota$ abstractly. $\Poly_{(-)}$ is a symmetric monoidal
  2-functor (SS Lem.\,4.4), so $\Poly_\E$ is a monoid in monoidal categories and colax functors. Eq.\,(16) records only
  the \emph{shape} of the map (and has a typo, $\tri$ for $\otimes$, in the second factor). The only explicit formula
  in SS is the Set one, Ex.\,4.6. The legs you use come from unwinding $\iota$ as the \emph{productor of}
  $\Poly_\times\colon\Poly_{\E\times\E}\to\Poly_\E$ (SS Prop.\,3.35, proof). That productor is the cartesian
  morphism obtained by applying $\times$ to the distributivity pullback and factoring through the canonical
  comparison. I checked this unwinding, and it gives exactly the legs you state. So nothing is wrong. But this is
  the step the identification actually rests on, and the note points the reader to the wrong place for it. Cite
  Prop.\,3.35 + Lem.\,4.4 and add two sentences of unwinding.
\item \textbf{(SHOULD) ``The single load-bearing coherence is $(\star)$'' overstates what is used.} The abstract
  (p.\,1), p.\,3 (``Its one coherence is $(\star)$''), Step 2 (``the only fact we verify by hand is that $(\star)$ is a
  genuine isomorphism'') and Table 1 row 3 all present the \emph{invertibility} of $\kappa$ as the proof's
  load-bearing input. But Step 3 uses $\kappa$ only as a \emph{map}, because it occurs definitionally on both
  sides. That is exactly what S-A made explicit. The invertibility of $\kappa$ is the cartesianness of the
  productor $\iota$, which SS prove (Prop.\,3.35: ``productor morphisms are cartesian''). Your proof neither
  re-proves it nor uses it. Honest framing: beyond the SS framework, the identification uses \emph{only chosen finite
  limits and universal properties}. $(\star)$ is where products meet pullbacks, but nothing about it is load-bearing.
  This also removes the slight tension with S-A(ii).
\item \textbf{(SHOULD) Size, and Rem.\,3.1.} The abstract says the result is ``not an artefact of \dots any size
  \dots hypothesis''. Rem.\,3.1 reduces along ``the fully faithful singular functor $\E\hookrightarrow\widehat\E$''.
  For large $\E$ (every Grothendieck topos), this is the Yoneda embedding into presheaves in a larger universe, as SS
  Ex.\,2.9 says explicitly. SS Lem.\,5.5 allows any small dense $F\colon\mathcal A\to\E$, and such an $F$ need not
  exist for a general $\E$. The SS inputs you cite (Thm.\,3.26, 5.6, Cor.\,5.12) are themselves proved through this
  reduction. So the honest statement is ``no size hypothesis beyond the universe enlargement SS's proofs already
  use''. Your own Steps 1--3 need no enlargement. Second point: a reduction along $F^*$ does not preserve ``on the
  nose''. $F^*$ of a chosen limit in $\E$ is \emph{a} limit in $\widehat{\mathcal A}$, not the chosen one. So Rem.\,3.1
  would give the identification only up to canonical isomorphism. Since the direct proof makes Rem.\,3.1
  unnecessary, I suggest cutting it to one sentence that carries both caveats, or deleting it.
\end{enumerate}

\textbf{NITs.}
(N1) Abstract, last line of the scope sentence: ``with $\Set$ as the degenerate case''. $\Set$ is the motivating
case. $\Vect$ is the degenerate one (your p.\,3).
(N2) S-A residues. Lemma 1.1 (p.\,2) still states the counit as $\varepsilon_c\otimes\varepsilon_d$, where \S3 has
$\lambda\circ(\varepsilon_c\otimes\varepsilon_d)$. And p.\,3 says ``the only isomorphism in sight being \dots
$\kappa$'', which omits $\lambda$.
(N3) \S4 (p.\,6): ``the category of $\tri$-comonoids is isomorphic to the category of internal categories and internal
cofunctors''. SS Cor.\,5.12 restricts to internal categories \emph{with exponentiable source maps}. Say so.
(N4) p.\,1: ``$\Cat^\sharp\ne\Cat$, as \S4 will show''. \S4 shows that $C\times D$ is not the product in
$\Cat^\sharp$. That does not by itself show that the categories are inequivalent. Rephrase as ``and the difference
matters, as \S4 shows''.
(N5) Table 1 (p.\,7) runs past the right margin in row 1, column 3.
(N6) Prior art. Spivak 2202.00534 \S8 already states in one sentence that $\Comon(\Poly)$ ``has finite sums and
products \dots\ it also has another monoidal product $\otimes$ with unit $y$, which agrees with the usual cartesian
product of categories''. Cite it next to NS Prop.\,8.79. It also bears on your question.

\section{Your question: which second result?}

I recommend the \textbf{Day-convolution identification}, stated as an $\E$-indexed (self-indexed) Day convolution
of the polynomial functors $P_p,P_q$ with respect to $\times$. It completes the note's own thesis and justifies its
history, since ``Day convolution'' is what the $\Set$ case was. It also looks short. $P_p=\Sigma_{A_p}y^{s_p}$ is an
indexed sum of representables. Co-Yoneda gives $y^{B}\circledast y^{B'}\cong y^{B\times B'}$. In each variable, Day
convolution should commute with indexed sums by Frobenius ($\Sigma_f\dashv f^*$ and stability of $\times$). So I expect
it to hold over any LCCC, with Grothendieck-ness needed at most to make the large coend exist. That expectation is
unchecked, and it is the main thing to test. The categorical product in $\Cat^\sharp(\E)$ is a different project.
Over $\Set$, existence is already in Spivak 2202.00534 \S8, via the cofree comonoid $\mathfrak c$. Internally it would
need cofree comonoids in $\Poly_\E$, so countable limits and exponentiability of infinite-depth tree bundles. That is
a paper of its own, not a second result for a 9-page note. Both directions are \emph{unchecked for novelty}. The
indexed-Day statement in particular may be folklore (Weber, Kock, or the indexed parts of 2202.00534), so run that
search before writing.

\end{document}
